% Solution: MWG (1995) Exercise 8.C.1 (IESDS and never-best-response elimination). Input by \solitem.
% Check: check_L2.py.
\begin{exsol}{MWG Exercise 8.C.1}{262}{Advanced}{3}
\asked Argue that if iterated elimination of strictly dominated strategies leaves a unique
prediction, iterated elimination of strategies that are never a best response leaves the same one.

\method A two-inclusions argument. A strictly dominated strategy is never a best reply, so
never-best-response elimination removes \emph{at least} what IESDS removes, round by round; and
elimination of never-best-replies can never remove everything, because a best reply exists in a
finite game. Definitions: iterated strict dominance (MWG \S8.B) and never-best-response elimination,
MWG Definitions~8.C.1--8.C.2 (rationalizable strategies; MWG \S8.C, pp.~242--245); the slides use the same idea for
rationalizability (Lecture~2, slides~11--13). The game is finite.

\solution
\textbf{Notation.} $D^0_i = R^0_i = S_i$. In round $k$,
\begin{itemize}
  \item $D^k_i$ = strategies in $D^{k-1}_i$ not strictly dominated (by a pure or mixed strategy
  over $D^{k-1}_i$) against every $s_{-i} \in D^{k-1}_{-i}$: iterated strict dominance;
  \item $R^k_i$ = strategies in $R^{k-1}_i$ that are a best reply, among all of $S_i$, to some belief over
  $R^{k-1}_{-i}$: iterated removal of never-best-replies, whose limit $R^\infty$ is the set of
  rationalizable strategies (MWG Definition~8.C.2).
\end{itemize}
Let $D^\infty = \bigcap_k D^k$ and $R^\infty = \bigcap_k R^k$; suppose $D^\infty = \{s^*\}$.

\textbf{Step 1: $R^k \subseteq D^k$ for every $k$.} By induction. True for $k = 0$. Suppose
$R^{k-1} \subseteq D^{k-1}$, and take $s_i \in R^k_i$ with $s_i \notin D^k_i$. Then $s_i$ is strictly
dominated by some $\sigma_i$ against every $s_{-i} \in D^{k-1}_{-i}$, hence against every belief
supported on $D^{k-1}_{-i}$ (an expectation of strictly positive differences is strictly positive). Beliefs over $R^{k-1}_{-i}$ are supported on
$R^{k-1}_{-i} \subseteq D^{k-1}_{-i}$, so against every belief over $R^{k-1}_{-i}$ the strategy $s_i$ is
strictly worse than $\sigma_i$ and hence never a best reply. That contradicts $s_i \in R^k_i$. So
$R^k_i \subseteq D^k_i$. Taking the limit, $R^\infty \subseteq D^\infty$.

\textbf{Step 2: $R^\infty \ne \emptyset$.} Each $R^k_i$ is non-empty by induction: given $R^{k-1}_{-i} \ne
\emptyset$, take any belief over it; a best reply in $S_i$ exists because $S_i$ is finite, and it
lies in $R^{k-1}_i$ (if it had been deleted in an earlier round $m \le k-1$, it would not be a best reply
to any belief over $R^{m-1}_{-i} \supseteq R^{k-1}_{-i}$). So it lies in $R^k_i$. The sets are finite,
nested and non-empty, hence their intersection $R^\infty$ is non-empty.

\textbf{Step 3: conclusion.} $\emptyset \ne R^\infty \subseteq D^\infty = \{s^*\}$ forces
$R^\infty = \{s^*\}$: elimination of never-best-replies yields the same unique prediction. $\blacksquare$

\answer{Every strictly dominated strategy is never a best reply, so round by round never-best-reply
elimination removes at least what IESDS removes: $R^k \subseteq D^k$. Since a best reply always exists,
$R^\infty$ is non-empty. If IESDS leaves only $s^*$, then $\emptyset \ne R^\infty \subseteq \{s^*\}$, so the
rationalizable set is also $\{s^*\}$.}

\interp Both procedures encode ``rational players who know that the others are rational''
(Lecture~2, slide~7 for IESDS, slides~11--13 for rationalizability), and the exercise says that the
second is never weaker. For two players the two sets are in fact equal (a strategy is never a
best reply iff it is strictly dominated, Pearce's lemma, noted in MWG \S8.C), so the content is
in games with three or more players with independent beliefs, where a strategy can be never a best reply
without being strictly dominated and rationalizability can be strictly sharper. Together with
exercise 5.1 this gives the chain
\begin{center}\small
strict dominant-strategy eq.\ $\Rightarrow$ unique IESDS survivor $\Rightarrow$
unique rationalizable profile $\Rightarrow$ unique Nash equilibrium.
\end{center}

\verify \texttt{check\_L2.py} computes both procedures on $500$ random finite two-player games (strategy
counts $2$--$4$, payoffs in $0,\dots,6$): IESDS with mixed dominators by linear programming, and
rationalizability by a feasibility LP over mixed beliefs. In every game $R^\infty \subseteq D^\infty$ and
$R^\infty$ is non-empty; whenever $D^\infty$ is a single profile (more than $30$ games) $R^\infty$ is the same
profile. (In two-player games they coincide, as the interpretation says.) MWG's solutions manual covers
this exercise; this argument is self-contained.

\mistakes
\begin{itemize}
  \item Claiming the two sets are equal in general: with three or more players and independent
  beliefs, $R^\infty$ can be strictly smaller than $D^\infty$.
  \item Forgetting Step~2: ``$R^\infty \subseteq \{s^*\}$'' does not give equality unless $R^\infty$ is
  non-empty, and non-emptiness is where finiteness (existence of a best reply) is used.
  \item Letting dominance be pure-only: with mixed dominators the set $D^\infty$ is the right one; with pure
  dominators the inclusion $R \subseteq D$ still holds, but $D$ is larger.
\end{itemize}
\end{exsol}
